Dear Editors of *Cognitive Systems Research*,

We respectfully submit our manuscript, **"The Regulatory Intelligence Program: Stability-First Cognitive Architectures Under Non-Learning Constraints,"** for consideration in *Cognitive Systems Research*.

This work presents a paradigm-level contribution to cognitive architecture research. Rather than treating intelligence as the optimization of external task performance, the paper advances **Regulatory Intelligence (RI)**—a framework in which intelligence is defined as a system's capacity to maintain internal viability (coherence, bounded dynamics, and recoverability) under cognitive stress.

The manuscript serves as an integrative overview of a multi-paper research program centered on *SpiralBrain v3.0*, a synthetic neurosymbolic architecture designed explicitly as a **non-learning, clean-slate cognitive system**. The goal of the program is not to compete with benchmark-optimized models, but to study cognition as a regulated dynamical process with observable internal physiology.

Key characteristics of the work that we believe align strongly with *Cognitive Systems Research* include:

* **Cognition as regulation:** The architecture treats cognition as homeostatic state evolution on a bounded 128-dimensional manifold, governed by explicit regulatory mechanisms rather than learned optimization.
* **Elastic (non-plastic) adaptation:** The system exhibits within-run regulatory adaptation while enforcing strict cross-run non-learning constraints, verified through explicit falsification tests.
* **Observable cognitive physiology:** Internal metrics such as coherence, drift, hazard, recovery time, and affective control signals are first-class experimental outputs, enabling analysis beyond task accuracy.
* **Affective control as computation:** A four-dimensional Symbolic-Emotional Calibration (SEC) signal functions as an upstream control mechanism for cognitive mode selection, rather than as an auxiliary annotation.
* **Architectural stability phenomena:** The work reports an empirically discovered phase-lock stability region in multi-pathway cognitive architectures, interpreted as a control-theoretic trade-off between differentiation and coherence.

Importantly, the manuscript is explicit about scope and non-claims. It does not propose a learning algorithm, a scaling strategy, or a benchmark-competitive model. Standard benchmarks are treated as **cognitive stressors**, not optimization targets, and lower task performance is interpreted as an intentional consequence of stability-preserving regulation rather than as a deficiency.

We believe this framing aligns closely with *Cognitive Systems Research*'s emphasis on cognitive architectures, dynamical systems, regulation, and interpretability. The contribution is intended to be read as a systems-level investigation into how stable cognition can be engineered and measured, rather than as a performance-driven AI study.

All experimental artifacts, logs, and analysis scripts referenced in the manuscript are available in a public repository to support transparency and independent inspection. The full cognitive system implementation is not publicly released; however, all reported results are derived from executable artifacts and logged outputs included for reproducibility of the reported findings.

Thank you for your consideration. We appreciate the journal's focus on foundational cognitive system design and believe this work will be of interest to readers concerned with stability, regulation, and the scientific study of cognition in artificial systems.

Sincerely,\\
John H. Cragin\\
Independent Researcher\\
[john.cragin@outlook.com](mailto:john.cragin@outlook.com)